% \section{Equações do Movimento Orbital e Manobras Associadas ao Problema de RVD/B}
% \label{sec_eq_mov_orbital_manobras}

% A formulação da dinâmica orbital no contexto de missões de RVD/B leva em consideração as diferentes fases da missão.
% Quando o lançador atinge a altitude nominal, é realizada a injeção do perseguidor em órbita coplanar à órbita do alvo, órbita inicial, a $100$~km abaixo da órbita do alvo, que atenda condições orbitais estáveis para então iniciar as manobras de longa, curta e de proximidade subsequentes.

% \subsubsection{Detalhes sobre a geometria da órbita do perseguidor}

% Para que o veículo lançador insira a espaçonave perseguidora na órbita de injeção desejada, é necessário fornecer os elementos orbitais \cite{vallado2007fundamentals} de injeção ou o vetor posição--velocidade no referencial inercial (ECI) .
% Neste trabalho foram considerados os elementos orbitais clássicos $(a, e, i, \Omega, \omega, M)$ e o vetor posição no referencial inercial no instante de injeção orbital.
% Além disso, as condições para que a injeção orbital seja considerada válida incluem:

% \begin{itemize}
%     \item o correto posicionamento da altitude do lançador na altitude desejada;
%     \item a fase orbital depende da janela de lançamento, podendo ser ajustada após a injeção e corrigida por espera, como no caso de deriva orbital.
% \end{itemize}


